\documentclass[11pt]{article}
\usepackage[T1]{fontenc}
\usepackage[utf8]{inputenc}
\usepackage{lmodern}
\usepackage[margin=1in]{geometry}
\usepackage{amsmath,amssymb,amsthm,mathtools,bm}
\usepackage{booktabs,graphicx,microtype}
\usepackage[numbers,sort&compress]{natbib}
\usepackage[colorlinks=true,allcolors=blue]{hyperref}
\usepackage[nameinlink,noabbrev]{cleveref}
\numberwithin{equation}{section}
\newtheorem{theorem}{Theorem}[section]
\newtheorem{proposition}[theorem]{Proposition}
\newtheorem{lemma}[theorem]{Lemma}
\newtheorem{corollary}[theorem]{Corollary}
\theoremstyle{definition}
\newtheorem{definition}[theorem]{Definition}
\newtheorem{assumption}[theorem]{Assumption}
\theoremstyle{remark}
\newtheorem{remark}[theorem]{Remark}
\newcommand{\R}{\mathbb R}
\newcommand{\E}{\mathbb E}
\newcommand{\Prob}{\mathbb P}
\newcommand{\dd}{\,\mathrm d}
\newcommand{\1}{\mathbf 1}
\newcommand{\pair}[2]{\langle #1,#2\rangle}
\newcommand{\norm}[1]{\left\lVert #1\right\rVert}
\newcommand{\abs}[1]{\left\lvert #1\right\rvert}
\DeclareMathOperator{\supp}{supp}
\title{Supplementary note: finite-grid preparation and initial-rate algebra
for additive coagulation--fragmentation}
\author{}
\date{}
\begin{document}
\maketitle

\section*{Setting}

This note accompanies the paper \emph{Sampling-law separation and finite
nonlinear corrections in additive coagulation--fragmentation} and uses its
notation. Sizes lie in $E=(0,\infty)$; $n_t$ is the finite nonnegative
measure of particle number per unit volume, $N(t)=n_t(E)$ its count, and
$M_1(t)=\int_E x\,n_t(\dd x)$ its mass, conserved at the value $m>0$.
Coagulation has a measurable symmetric kernel $K(x,y)\geq0$, and
fragmentation has a per-particle selection rate $S(x)\geq0$ with an
expected daughter measure $b_x$ on $(0,x)$ of total count two and total
mass $x$. The additive model of the main paper is $K(x,y)=\lambda(x+y)$
and $S\equiv\sigma$, with $b=\lambda m$. For a finitely supported
preparation $v$, its count vector field is
\[
 \mathcal R(v)=\int S(x)v(\dd x)-\frac12\iint K(x,y)v(\dd x)v(\dd y).
\]
The count-neutral classification theorem of the main paper states that
$\mathcal R(v)=0$ for every nonnegative finitely supported $v$ with
$M_1(v)=m$ if and only if there is a nonnegative pointwise finite
measurable function $d$ with
\begin{equation}
 K(x,y)=xd(y)+yd(x),\qquad S(x)=md(x).
 \tag{$\ast$}\label{eq:neutral-class-supp}
\end{equation}

\section{Preparation constraints and elementary initial-rate tomography}

This note records supporting algebra for the observation limitations
discussed in the main paper's section on bulk observables. Quadratic response models on mixture simplices
are classical; see, for example, \citet{BrownDonevBissett2015}.
The statements below are direct applications of quadratic algebra,
not a separate inverse-problem method. They concern exact initial rates
under specified preparations. Full-trajectory conclusions are stated
only when a conserved quantity or a closed count equation justifies them.

\subsection{Quadratic rates on a finite preparation shell}

On a grid of $q$ selected sizes, let $z\in\R^q$ be the vector of number
concentrations, $\mathsf K$ a symmetric kernel matrix, and $s$ the vector
of single-particle count production rates. Write
$\mathcal R(z)=s^Tz-z^T\mathsf Kz/2$. The entries need not be positive
for the following algebraic statement; physical admissibility is an
additional constraint.

\begin{proposition}[Finite preparation-shell classification]
\label{prop:preparation-shell}
Let $\mathsf C\in\R^{r\times q}$ have full row rank, let
$h\in\R^r\setminus\{0\}$, and suppose
$\{z>0:\mathsf Cz=h\}$ is nonempty. Then $\mathcal R$ vanishes on this
shell if and only if there is $\mathsf A\in\R^{r\times q}$ such that
\begin{equation}
 \mathsf K=\mathsf C^T\mathsf A+\mathsf A^T\mathsf C,
 \qquad s=\mathsf A^Th.
 \label{eq:preparation-shell}
\end{equation}
In particular $\operatorname{rank}\mathsf K\leq2r$.
\end{proposition}

\begin{proof}
A polynomial vanishing on a nonempty relatively open subset of an affine
space vanishes on the whole affine space. Put
$P=\mathsf C^T(\mathsf C\mathsf C^T)^{-1}\mathsf C$ and $Q=I-P$.
The Hessian restricted to $\ker\mathsf C$ vanishes, so polarization gives
$Q\mathsf KQ=0$. Define
\[
 \mathsf A_0=(\mathsf C\mathsf C^T)^{-1}\mathsf C\mathsf K(I-P/2).
\]
A direct expansion using $Q\mathsf KQ=0$ gives
$\mathsf C^T\mathsf A_0+\mathsf A_0^T\mathsf C=\mathsf K$.
On the affine space the residual is
$(s-\mathsf A_0^Th)^Tz$. Its vanishing implies
$s-\mathsf A_0^Th=\mathsf C^T\gamma$ and $\gamma^Th=0$:
the residual annihilates $\ker\mathsf C$, and evaluation at one feasible
$z$ gives the second identity. The skew matrix
\[
 L=\frac{h\gamma^T-\gamma h^T}{h^Th}
\]
satisfies $L^Th=\gamma$. Hence $\mathsf A=\mathsf A_0+L\mathsf C$
keeps the representation of $\mathsf K$ and gives $\mathsf A^Th=s$.
The converse follows by substitution, and the rank bound follows from
\eqref{eq:preparation-shell}.
\end{proof}

The rank bound is necessary, not sufficient for invisibility.
For simultaneous constraints on count and mass, the rows of
$\mathsf C$ are $(1,\ldots,1)$ and the grid sizes, with $h=(c,m)^T$.
A corresponding continuum sufficient family is
\begin{equation}
 K(x,y)=a(x)+a(y)+xd(y)+yd(x),\qquad S(x)=ca(x)+md(x).
 \label{eq:two-constraint-family}
\end{equation}
For bounded nonnegative $a,d$, every mass-conserving solution with a
valid count balance has
\[
 N'=(c-N)\int a\dd n_t+(m-M_1)\int d\dd n_t.
\]
If $N(0)=c$ and $M_1=m$, the locally integrable coefficient
$\int a\dd n_t$ and the scalar linear equation give $N(t)=c$.
This is a conditional statement about such solutions, not a new forward
existence theorem. The finite-grid classification does not by itself
prove a continuum converse under two constraints. In contrast,
the count-neutral classification theorem of the main paper proves the
complete continuum classification when only initial mass is fixed and
count is allowed to vary.

\subsection{Dilution, sources, and the exact algebraic ambiguity}

For this subsection allow a size-independent count source $J$ and a net
single-particle count production rate $\beta(x)$; these are observation
parameters, distinct from $b=\lambda m$ in the additive model.
For a preparation $n_0=cp$, where $p$ is a finitely supported probability
measure, suppose the observed initial count rate is
\begin{equation}
 \mathfrak r(c,p)=J+c\int\beta\dd p
                 -\frac{c^2}{2}\iint K(x,y)p(\dd x)p(\dd y).
 \label{eq:initial-rate}
\end{equation}
Assume symmetric, pointwise finite $K$, known positive concentrations,
and unchanged $J,\beta,K$ across preparations. Antisymmetric parts of a
kernel never enter this observation, which explains the symmetry assumption.

\begin{proposition}[Fixed-concentration gauge]\label{prop:dilution-gauge}
With known $J$, two models agree in \eqref{eq:initial-rate} for all $p$
at one fixed concentration $c$ if and only if their differences satisfy
\begin{equation}
 \Delta K(x,y)=a(x)+a(y),\qquad\Delta\beta(x)=ca(x)
 \label{eq:known-source-gauge}
\end{equation}
for some function $a$. Agreement at two distinct positive concentrations
removes this ambiguity. With unknown $J$, agreement at two distinct
positive concentrations $c_1,c_2$ leaves exactly the algebraic family
\begin{equation}
 \Delta J=\zeta,\qquad
 \Delta\beta=-\zeta(c_1^{-1}+c_2^{-1}),\qquad
 \Delta K=-\frac{2\zeta}{c_1c_2}.
 \label{eq:unknown-source-gauge}
\end{equation}
A third distinct concentration removes this family.
\end{proposition}

\begin{proof}
For known $J$, monodisperse preparations give
$\Delta\beta(x)=c\Delta K(x,x)/2$. Equal mixtures of two atoms give
$\Delta K(x,y)=[\Delta K(x,x)+\Delta K(y,y)]/2$.
These are \eqref{eq:known-source-gauge}, and substitution proves
sufficiency. Applying the diagonal identity at two concentrations
forces both differences to vanish.
For unknown $J$, monodisperse data at $c_1,c_2$ give the two linear
equations
$0=\zeta+c_j\Delta\beta(x)-c_j^2\Delta K(x,x)/2$.
Solving gives the constants in \eqref{eq:unknown-source-gauge}; equal
mixtures then give the same constant off the diagonal. Their rate
difference at any concentration $c$ is
$\zeta(1-c/c_1)(1-c/c_2)$, proving both remaining assertions.
\end{proof}

Nonnegativity can restrict the admissible values of $\zeta$ or force
$\zeta=0$ at a boundary model. Thus the algebraic family does not assert
physical nonuniqueness at every parameter value. Diluting the same shape
$p$ changes both count and material loading. It differs from changing
$p$ while preserving $\int x\dd n_0=m$, which leaves the entire count-neutral
class \eqref{eq:neutral-class-supp} invisible.

\subsection{Finite-grid reconstruction and noise bounds}

For known $J$ and any $c>0$, direct substitution gives
\begin{align}
 K(x,y)&=\frac{\mathfrak r(c,\delta_x)+\mathfrak r(c,\delta_y)
 -\mathfrak r(2c,(\delta_x+\delta_y)/2)-J}{c^2},
 \label{eq:kernel-tomography}\\
 \beta(x)&=\frac{4\mathfrak r(c,\delta_x)
                 -\mathfrak r(2c,\delta_x)-3J}{2c}.
 \label{eq:linear-tomography}
\end{align}
On $q$ selected sizes this uses $q$ monodisperse rates at $c$ and
$q(q+1)/2$ equal-mixture rates at $2c$, including repeated sizes.
The total equals the number of unrestricted symmetric-kernel and
linear-rate parameters. The observation map is linear in those
parameters; fewer scalar initial-rate observations cannot be injective
on an open parameter set. This is only a measurement-count bound.
It makes no claim of optimal noise performance or optimal use of
finite-time trajectories. If $J$ is unknown, one further rate suffices,
because for a reference size $x$,
\begin{equation}
 J=3\mathfrak r(c,\delta_x)-3\mathfrak r(2c,\delta_x)
                  +\mathfrak r(3c,\delta_x).
 \label{eq:source-tomography}
\end{equation}

With exact $J$ and absolute errors at most $\varepsilon$ in each rate,
\cref{eq:kernel-tomography,eq:linear-tomography} imply
\begin{equation}
 |\widehat K-K|\leq3\varepsilon/c^2,\qquad
 |\widehat\beta-\beta|\leq5\varepsilon/(2c).
 \label{eq:rate-noise}
\end{equation}
If $J$ is estimated with error at most $\varepsilon_J$, the right sides
increase by $\varepsilon_J/c^2$ and $3\varepsilon_J/(2c)$, respectively.
Using \eqref{eq:source-tomography} with the same rate-error bound gives
$\varepsilon_J\leq7\varepsilon$; correlations among reused observations
do not affect these deterministic inequalities.

Suppose now the count difference $N(t)-N(0)$ has error at most $\eta$
and each required trajectory has $|N''(s)|\leq L$ for $0\leq s\leq t$.
Taylor's integral remainder gives initial-rate error at most
$\eta/t+Lt/2$. With exact $J$ this yields
\begin{equation}
 |\widehat K-K|\leq\frac3{c^2}\left(\frac\eta t+\frac{Lt}2\right).
 \label{eq:short-time-noise}
\end{equation}
When $\eta,L>0$, the bound is minimized at $t=\sqrt{2\eta/L}$ if
that time lies in the interval where the curvature bound is known;
the resulting bound is $3\sqrt{2L\eta}/c^2$. Zero-error or zero-curvature
cases are read directly from \eqref{eq:short-time-noise}.
An exact initial count and terminal count error at most $\eta$ are
one way to meet the assumed count-difference error bound; if both are
noisy, their errors must be combined.

A uniform curvature bound is an additional assumption and can worsen
with concentration, so these formulas do not justify arbitrarily large
$c$. Finite-width preparations measure kernel averages; localization
error requires a separate regularity bound. Dilution must preserve the
kinetic coefficients for \eqref{eq:initial-rate} to apply across
experiments. Changes in the surrounding medium can violate this
assumption. Collision-induced fragmentation also contributes a quadratic
term and cannot be separated from coagulation by concentration scaling
alone. None of these count-rate observations identifies the daughter
measure; the distributional data used for late-time Fourier identification in
the main paper contain information that count measurements lack.

\bibliographystyle{plainnat}
\bibliography{../references}
\end{document}
